% ============================================================================
% When Safe Search Budgets Do Not Transfer Across Graph-ANNS Rebuilds
% Anonymous ICLR submission — revised manuscript (v2, P0–P7 evidence integrated)
% ----------------------------------------------------------------------------
% Compile with the official ICLR kit: place iclr2027_conference.sty,
% iclr2027_conference.bst (and the ICLR banner files, e.g. iclr2027_logo.pdf)
% next to this file, then run:  pdflatex -> bibtex -> pdflatex x2
% If the ICLR style kit is temporarily unavailable, uncomment the FALLBACK
% block below to compile with plain article class.
% Figures required (same directory): fig1_setting.png, fig2_workflow.png,
% fig3_crossfamily.png, fig4_pooling.png, fig5_plane.png, fig6_ablation.png,
% fig7_ladder.png  (shipped in the artifact media/ folder).
% ============================================================================
\documentclass{article}

% ---- ICLR official style if present; otherwise self-contained fallback ----
\IfFileExists{iclr2027_conference.sty}{%
  \usepackage{iclr2027_conference,times}
}{%
  \usepackage[margin=1in]{geometry}
  \usepackage{times}
  \usepackage{natbib}
  \usepackage{fancyhdr}
  \pagestyle{fancy}\fancyhf{}
  \fancyhead[R]{\small\itshape Safe Search Budgets Across Graph-ANNS Rebuilds}
  \fancyfoot[C]{\small Under review as a conference paper at ICLR 2027 --- page \thepage}
  \renewcommand{\headrulewidth}{0pt}
}

\usepackage[utf8]{inputenc}
\usepackage[T1]{fontenc}
\usepackage{amsmath,amssymb,amsthm,mathtools}
\usepackage{booktabs}
\usepackage{graphicx}
\usepackage{multirow}
\usepackage{tabularx}
\usepackage{microtype}
\usepackage[hidelinks]{hyperref}

\newtheorem{theorem}{Theorem}
\newtheorem{proposition}[theorem]{Proposition}
\newtheorem{corollary}[theorem]{Corollary}

\newcommand{\DistComp}{\mathrm{DistComp}}
\newcommand{\TV}{\mathrm{TV}}
\newcommand{\KL}{\mathrm{KL}}
\renewcommand{\Pr}{\mathop{\mathrm{Pr}}}   % \Pr is a kernel command: renew, not new
\newcommand{\E}{\mathbb{E}}
\newcommand{\Acal}{\mathcal{A}}
\newcommand{\Ecal}{\mathcal{E}}
\newcommand{\BOT}{\bot}

\title{When Safe Search Budgets Do Not Transfer Across Graph-ANNS Rebuilds}

\author{Anonymous Authors\\
Paper under double-blind review\\
\texttt{anonymized artifact accompanies the submission}}

\begin{document}

\maketitle

\begin{abstract}
Graph approximate nearest-neighbor systems rebuild indexes after data refreshes, maintenance, replica
creation, and randomized construction, yet downstream search-budget policies are commonly reused across
builds. We study when this transport is statistically justified. Each serialized build is modeled as an
algorithmic environment with an implementation-native finite action set, an endpoint-aware failure event,
and a categorical unresolved state. A transcript-limited two-environment reduction shows that, when an
available probe transcript cannot distinguish builds requiring conflicting useful actions, any build-blind
policy must incur unsafe execution, conservative computation, fallback, or abstention in at least one
environment; we instantiate this premise by measuring that registered build pairs have near-chance
hit-count transcripts at chance (plug-in total variation at most $0.25$) while their per-query responses conflict on
45--57\% of queries. A complementary fixed-target theorem gives sufficient conditions for safe recovery, and
we measure where those conditions hold: on the registered grid the positive-margin condition fails for all
48 targets, so certified recovery collapses to the maximum budget or abstention. Operationally, ICBA, a
certified portability audit, and a deterministic-rebuild contract convert these findings into a decision
rule. Across SIFT-100K and Arxiv-Nomic-100K, strictly comparable hnswlib and Faiss HNSW experiments exhibit
75.87--91.07\% finite-action safe-budget variation and 17.17--23.60\% incremental transport risk, and the
phenomenon replicates at SIFT-1M (21.87\%, 95\% CI $[21.52, 22.21]$). A max-over-source pooling baseline over
22 registered builds cuts risk below 1\% at $1.34$--$1.45\times$ per-query-oracle cost, though its source
requirement grows with scale; a learned query-level predictor fails to transfer across builds (risk 21--24\%,
transfer $R^2 = 0.015$). Fixing input order and single-threaded construction eliminates measured transport
risk with $1.71$--$3.52\times$ build overhead; order pinning alone halves response variation at no build-time
cost. Primitive profiling remains cheap at both scales, so we claim neither universal profiling cost nor
open-world recovery.
\end{abstract}

\section{Introduction}
\label{sec:intro}

Graph-based approximate nearest-neighbor search (Graph-ANNS) is often tuned as though an index were a
stable substrate. After rebuilding with unchanged data and nominal hyperparameters, practitioners may reuse
a global budget, query-level predictor, or stopping rule selected on an earlier index. The statistical
premise behind this practice is rarely stated: the policy is assumed to remain safe and useful under a new
realization of the construction algorithm. Reconfiguration can violate that premise through input order,
randomized state, parallel scheduling, graph pruning, software version, or build configuration---even when
the query and data distributions do not change.

The reconfiguration we study arises in ordinary operation. A vector collection is frequently built more
than once: serving replicas are (re)created on different machines with different insertion orders and
parallel schedules; autoscaling adds replicas; failure recovery and blue-green deploys rebuild from the same
data; library upgrades rebuild under new toolchains. A budget policy calibrated on one replica---a
canary---is then reused across the replica set. We study the risk of this plausible deployment shortcut:
whether calibrated budgets remain safe on independently built replicas of the same collection. We make no
claim that every operator reuses budgets and cite no production incident; the setting is stated as a
plausible, common-shape deployment pattern, and all conclusions are conditional on registered builds.

This issue is not simply a drop in mean recall. Two builds can have similar aggregate quality while
assigning different minimum safe actions to the same query. Reusing a source action can then under-budget a
target build; choosing a uniformly larger action can avoid some failures while paying conservative
computation. Existing adaptive search, cost prediction, and risk-control methods are valuable for a
profiled target index, but their calibration does not automatically transport across an independently
rebuilt index. Conversely, target profiling is not uniformly expensive: our primitive measurements are
sub-second at 100K and 1M scale. The scientific gap is the validity and decision semantics of transport, not
a blanket claim that profiling is prohibitive.

We formulate build reconfiguration as a finite environment--action problem. For each serialized build and
query, the safe budget is the smallest registered implementation-native action satisfying a quality event.
If no registered action succeeds, we retain a categorical unresolved symbol rather than imputing a larger
numerical budget. Directed transport separates four outcomes that are frequently conflated: unsafe
execution, conservative computation, certified fallback, and abstention. This separation turns an informal
robustness concern into an auditable risk--cost object.

Our evidence addresses three layers with different strength. First, the query-level safe-budget
\emph{response} is strongly non-portable across registered rebuilds (17.17--23.60\% incremental transport
risk at 100K, replicated at 1M). Second, at the \emph{global policy} layer we give a constructive answer:
pooling 22 source builds into a per-query conservative action cuts risk below 1\%, and certified fixed-action
fallback is possible but only at the maximum registered budget. Third, whether a \emph{learned query-level
predictor} trained on one build transfers to another is probed directly: it does not
(Section~\ref{sec:scale}), completing a measured three-layer boundary rather than an extrapolation.

Our theoretical contribution separates impossibility from conditional recovery. The negative result is
transcript limited and now empirically instantiated: registered build pairs are nearly indistinguishable at
the transcript level (plug-in total variation at most 0.25 over all 552 pairs at every probe level we
measure) while requiring conflicting actions on up to 57\% of queries. The positive result is fixed-target
and finite: endpoint feasibility, action-relevant information, a positive risk margin, independent target
evidence, and a valid fallback or abstention rule suffice for safe screening and certification---and our
per-target audit shows exactly which condition fails in practice (margin, on all 48 registered targets).
The testing inequalities are classical; the contribution is their build-conditioned Graph-ANNS
formulation, operational composition, and measured boundary.

We instantiate the theory in ICBA, a certified audit rather than a new traversal algorithm. ICBA validates
action semantics and role separation, measures cross-build response heterogeneity, evaluates directed
transport, and certifies only frozen fixed-target actions. The audit yields a cost-ordered, four-way
decision rule: enforce and replay-check a deterministic build contract when feasible; otherwise pool many
source builds when operational history permits; otherwise collect target evidence and certify a frozen
action; if none yields a certified action, abstain. A certificate cannot create a useful action when the
candidate family contains none.

The main empirical evidence comprises four strictly comparable HNSW data--implementation cells (hnswlib and
Faiss on SIFT-100K and Arxiv-Nomic-100K), a preregistered SIFT-1M replication cell, and two Vamana-style
supporting cells. Finite-action safe-budget variation is 75.87--91.07\% at 100K (98.4\% at 1M), and
incremental transport risk is 17.17--23.60\% with all confidence intervals above the preregistered 2\%
materiality threshold. Deterministic single-thread rebuilding reduces measured transport risk to zero in
registered controls while increasing build time by $1.71$--$3.52\times$; a descriptive chain ablation
attributes the entire guarantee to single-threading, showing order pinning alone halves variation for free
and seed pinning under parallel construction is inert.

Our contributions are:
\begin{itemize}\itemsep1pt
\item We define build-reconfiguration portability using implementation-native actions, categorical
unresolved endpoints, and four distinct deployment outcomes, with formal estimands for incremental
transport risk and the variation families.
\item We give a transcript-limited lower bound---whose near-indistinguishability premise we instantiate
empirically---and a fixed-target conditional recovery theorem whose five conditions we audit per target,
finding the margin condition systematically absent.
\item We introduce ICBA as a certified diagnostic workflow and derive a deterministic-rebuild contract whose
mechanistic ablation isolates the load-bearing clause (single-threading).
\item We provide clean cross-implementation HNSW evidence at 100K and a preregistered 1M replication,
constructive baselines (robust source pooling with measured scale boundaries; select-then-certify with its
measured failure), a learned-predictor transfer probe, and supporting Vamana-style evidence, with estimand,
uncertainty, and comparability boundaries stated explicitly.
\item We organize extensive negative method evidence as a theory-guided falsification ladder, identifying
whether each route fails at attainability, observability, certification, tail behavior, or economics.
\end{itemize}

\begin{table}[t]
\centering
\caption{Takeaway: deployment scenario $\rightarrow$ recommended route $\rightarrow$ measured boundary (all numbers registered-family conditional).}
\label{tab:takeaway}
\footnotesize
\setlength{\tabcolsep}{3pt}
\begin{tabular}{p{3.2cm}p{3.9cm}p{6.5cm}}
\toprule
Scenario & Route & Measured boundary \\
\midrule
Rebuild feasible offline; near-zero risk required & Deterministic contract (canonical order + seed + single thread) & risk 0; build $1.71$--$3.52\times$; order pinning alone: variation halved, risk not zero \\
$\ge$10 prior builds with per-query records & Robust source pooling (max or $q{=}0.9$) & 100K: risk $<1.1\%$ @ $1.34$--$1.45\times$ oracle; 1M: $k{=}7$ leaves 10.4\% (scale-dependent) \\
Cost-sensitive, risk tolerance moderate & Canonical order pinning only & free ($0.80$--$0.93\times$ build); residual variation 30--34\% \\
Fresh target, no history, budget for truth & Certify a frozen action & only max budget clears ($P{=}0.47$--$0.63$ at $m{=}59$); margin-band actions fail $m{=}94$ \\
No route yields a certified action & Abstain & safety preserved; regret bounded in Table~\ref{tab:constructive} \\
\bottomrule
\end{tabular}
\end{table}

Figure~\ref{fig:setting} summarizes the setting. Section~\ref{sec:related} positions the problem.
Sections~\ref{sec:formulation}--\ref{sec:theory} define the model and theory;
Section~\ref{sec:icba} presents ICBA and the mitigation contract;
Sections~\ref{sec:protocol}--\ref{sec:ladder} give the protocol, evidence, and falsification ladder;
Sections~\ref{sec:limitations}--\ref{sec:conclusion} state limitations and conclusions.

\begin{figure}[t]
\centering
\includegraphics[width=.8\linewidth]{fig1_setting.png}
\caption{Build reconfiguration induces a policy-transport problem. A source-calibrated action may
under-budget the target build, while a uniformly large action may be safe but conservative. The build is
hidden only relative to the available probe transcript; ICBA separates measurement, certification, and
value.}
\label{fig:setting}
\end{figure}

\section{Related Work}
\label{sec:related}

\paragraph{Graph-based ANNS construction and search.}
HNSW uses a randomized multilayer proximity graph and greedy candidate-list search
\citep{malkov2020hnsw}. NSG approximates a monotonic relative-neighborhood structure \citep{fu2019nsg}, and
DiskANN introduced Vamana-style robust pruning for memory- and disk-efficient search at scale
\citep{subramanya2019diskann}. Faiss provides an independent optimized HNSW implementation
\citep{johnson2021faiss}. Incremental graph maintenance can avoid full rebuilding after updates
\citep{xu2022maintenance}, and recent work revisits proximity-graph construction cost and pruning from
practice to theory \citep{prokhorenkova2020graph}. Insertion order and intrinsic dimensionality measurably
affect HNSW recall \citep{elliott2024order}, and per-index budget autotuning is an active systems topic
\citep{bae2026qbat}. These systems reduce construction or maintenance cost; our question is complementary:
when a build does change, does a previously selected safety-critical search action transport to the
realized graph?

\paragraph{Query-adaptive search and cost prediction.}
Learned early termination, adaptive exploration, conformal retrieval, and learned query-cost estimators
improve efficiency for a concrete profiled index
\citep{li2020learned,chatzakis2025darth,horchidan2025conann,wang2026annie}; distribution-aware adaptive
HNSW exploration extends the family \citep{zhang2026exploration}. They establish that query difficulty is
predictable within an environment. Unlike these methods, we do not propose another predictor; we test
whether predictor or budget selected on one serialized build retains its risk and cost semantics after an
independent rebuild---and our probe shows a representative regressor trained on 23 source builds transfers
at chance level (Section~\ref{sec:scale}).

\paragraph{Risk control and safe selection.}
Learn-Then-Test, conformal risk control, selective classification, and safe best-arm identification provide
finite-sample tools for selecting or abstaining within a fixed candidate family
\citep{angelopoulos2025ltt,bates2021rcps,angelopoulos2024crc,elyaniv2010selective,wang2022safebai}; the
fixed-confidence sample complexity of such selections parallels best-arm identification
\citep{garivier2016bai}. ICBA uses these classical tools only for a named target and frozen actions. Its
distinct contribution is to expose when build-conditioned transport invalidates the premise of reuse, and
to keep unsafe execution, conservative cost, fallback, and abstention as separate decisions.

\paragraph{Distribution shift and algorithmic reconfiguration.}
Domain adaptation theory shows that source performance need not determine target performance without
support or structural assumptions \citep{bendavid2010theory,johansson2019support}. Multi-environment
predictive inference treats environments as higher-level sampling units \citep{duchi2025multienv}.
Rebuilding an index is neither ordinary covariate shift nor a change in the underlying vector dataset: the
query distribution may remain fixed while the algorithmic state changes. We therefore treat each serialized
build and its replay semantics as an environment. This framing supports a precise fixed-target guarantee
while making clear that query bootstrap conditional on registered builds is not an unseen-build confidence
interval.

\section{Problem Formulation}
\label{sec:formulation}

Let $\Ecal$ be a preregistered finite set of build environments. An environment $E\in\Ecal$ consists of a
serialized index together with its implementation, construction configuration, input permutation, seed when
meaningful, toolchain, and replay semantics. Each implementation has its own finite native action set
$\Acal_E$, and queries follow a frozen law $q\sim P_Q$.

For action $a\in\Acal_E$, define the endpoint-aware failure indicator and implementation-internal cost by
\begin{equation}
Z_{E}(q,a)=\mathbf{1}\{\text{the returned result fails the registered quality target}\},\qquad
C_{E}(q,a)\ge 0,
\label{eq:zc}
\end{equation}
and the minimum observed safe budget by
\begin{equation}
B_{E}(q)=\min\{a\in\Acal_{E}:Z_{E}(q,a)=0\},\qquad
B_{E}(q)=\BOT\ \text{if no registered action succeeds}.
\label{eq:bmin}
\end{equation}
The symbol $\BOT$ is categorical and is never imputed as the largest action or a larger numerical budget. A
policy observes a transcript $T$ and query-side observables, then returns an action or abstains. Its
absolute target risk is
\begin{equation}
r_{E}^{\pi}=\Pr_{q\sim P_{Q}}\!\left[Z_{E}(q,\pi_{T,q})=1\right].
\label{eq:risk}
\end{equation}
We use risk tolerance $\delta=0.05$, certification error probability $\alpha=0.05$, and safety margin
$\gamma>0$ stated per experiment; the materiality gate $0.02$ is a preregistered scientific threshold---chosen, before any 100K evaluation
role was accessed, as the smallest risk level that would still be operationally material for budget
reuse---not a utility function. ``Preregistered'' in this paper means frozen in a versioned manifest with commit hash and
timestamp before the corresponding evaluation role was accessed; it does not refer to an external OSF
registration. A rejection event is $R$. An action $a_f$ is a fallback only if it has independent safety
evidence under the same target and failure event; otherwise rejection produces abstention.

The primary estimand of Section~\ref{sec:results} is the stage-registered endpoint-aware incremental
transport risk
\begin{equation}
\Delta_{s\rightarrow t}
=\Pr_{q\sim P_{Q}}\!\left[Z_{t}\big(q,B_{s}(q)\big){=}1\right]-r_{t}^{\mathrm{ref}},
\label{eq:delta}
\end{equation}
where the reference risk on the target follows the stage-registered reference construction; for the four
HNSW cells this is the reference event of the repaired $h{=}10$ common estimand. Intervals are 95\%
query-bootstrap over 5{,}000 resamples (seed 991) retaining each sampled query's full vector of directed-pair
observations, conditional on the registered build family. Heterogeneity is reported through two separated
variation families,
\begin{equation}
V_{\mathrm{fin}}=\Pr_{q}\!\left[\big|\{B_{E}(q):E\in\Ecal_{\mathrm{reg}},\,B_{E}(q)\neq\BOT\}\big|>1\right],
\qquad
V_{\mathrm{end}}=\Pr_{q}\!\left[\mathbf{1}\{B_{E}(q){=}\BOT\}\ \text{nonconstant over }E\right],
\label{eq:var}
\end{equation}
finite-action variation and endpoint-state variation; reported separately so response heterogeneity cannot
be attributed to unresolved endpoints. We distinguish registered-grid unresolved, true endpoint infeasible,
and right-censored queries. Current finite grids often cannot separate the latter two, so we report
unresolved mass without promoting it to true infeasibility.

\begin{table}[t]
\centering
\caption{Implementation-native semantics. Raw action values and distance-computation counts are never
compared across implementations.}
\label{tab:semantics}
\footnotesize
\setlength{\tabcolsep}{4pt}
\begin{tabular}{p{2.4cm}p{3.4cm}p{2.1cm}p{5.2cm}}
\toprule
Implementation & Registered build environment & Native action & Semantic guardrail \\
\midrule
hnswlib HNSW & seed $\times$ input order (24 builds/dataset) & \texttt{efSearch} & candidate-list control; not a DC cap \\
Faiss HNSW & registered input permutation (24 builds/dataset) & \texttt{efSearch} & implementation-specific; not equated \\
DiskANN3/Vamana-style & registered permutation; fixed options (12/dataset) & $l$, beam width $=1$ & list size; $\neq$ HNSW \texttt{efSearch} \\
\bottomrule
\end{tabular}
\end{table}

\paragraph{Cost naming.}
ROM-NDC quantities in earlier drafts are renamed \emph{DistComp} throughout:
$\DistComp(x \text{ vs } y)$ is the ratio of mean numbers of distance computations of $x$ to $y$, computed
only within one implementation. ``NDC'' alone is avoided to prevent confusion with NDCG.

\section{Theory}
\label{sec:theory}

\subsection{Main Theorem I: transcript-limited non-portability}

Consider two build environments $E_0,E_1$. Under $E_i$, a procedure observes $T\sim P_i^{T}$ and outputs an
action through any randomized Markov kernel measurable with respect to $T$. Let $D_i$ be the region of
nontrivially acceptable actions in $E_i$, assume $D_0\cap D_1=\varnothing$, and let
$l_i\ge\Delta\,\mathbf{1}\{a\notin D_i\}$ for some $\Delta>0$.

\begin{theorem}[transcript-limited lower bound]
Every transcript-based randomized procedure satisfies
\begin{equation}
\max_{i\in\{0,1\}}\E[l_{i}]\ \ge\ \frac{\Delta}{2}\bigl(1-\TV(P_{0}^{T},P_{1}^{T})\bigr).
\label{eq:thm1}
\end{equation}
\end{theorem}

The decision regions encode actions that avoid conservative cost in one environment and actions that avoid
under-budget failure in the other. If the transcript available under the registered probe budget cannot
identify which region applies, no transcript-measurable policy can uniformly avoid unsafe execution,
conservative computation, fallback, or abstention. The statement quantifies over randomized policies but is
conditional on the transcript laws; it neither claims that every build pair is indistinguishable nor
forbids target probing.

For an adaptive transcript with $m$ observations and KL divergence $K_m$ between the transcript laws, the
classical Bretagnolle--Huber and Pinsker inequalities give
\begin{equation}
\max_{i}\E[l_{i}]\ \ge\ \frac{\Delta}{4}e^{-K_{m}},\qquad
\max_{i}\E[l_{i}]\ \ge\ \frac{\Delta}{2}\max\!\Bigl(0,\,1-\sqrt{K_{m}/2}\Bigr).
\label{eq:bh}
\end{equation}
Le Cam reduction, total variation, KL chain rules, and the testing inequalities are classical
\citep{yu1997assouad}; the domain-specific contribution is the reduction from a wrong build decision to the
four Graph-ANNS deployment outcomes.

\paragraph{Empirical instantiation (Section~\ref{sec:scale}).}
We measure this premise directly, for two probe classes. Over all 552 registered build pairs and probe
levels $k\in\{1,3,6\}$ cheapest actions, a strong classifier attributes \emph{hit-count} transcripts to
their build at chance accuracy (median 0.49--0.51), and the exact plug-in total variation between
transcript marginals never exceeds $0.25$---while the same pairs disagree on the per-query minimum safe
action for 45--57\% of queries. For \emph{runtime} probes (visited nodes, queue statistics, median
latency at the cheapest action), separation is weak-to-moderate: median accuracy 0.56--0.58, worst pair
0.72, i.e., TV $\ge 2(0.72)-1=0.44$ on the most separable pair. Near-indistinguishability therefore
holds for the hit-count transcript class and only partially for runtime features. Substituting the
worst-case $\TV=0.44$ into Theorem~1 still bounds the unavoidable loss below by $0.28\,\Delta$ for such
pairs: the deployment-loss conclusion survives quantitatively, the ``indistinguishable'' wording is
probe-class-qualified, exactly as the theorem requires.

\subsection{Main Theorem II: finite-class conditional recoverability}

Fix a named target $E_t$, a frozen family $\Acal=\{a_1,\dots,a_M\}$, a declared $P_Q$, and the
endpoint-aware loss. Suppose: (i) at least one action has risk at most $\delta-2\gamma$; (ii) the selection
transcript distinguishes action-relevant alternatives in the registered class; (iii) a useful action has
positive risk margin; (iv) selection and final certification use independent query roles and the candidate
family is frozen before certification; and (v) rejection invokes an independently justified fallback or
abstention. Let the simultaneous estimation event be
\begin{equation}
\mathcal{G}
=\Bigl\{\sup_{a\in\Acal}|\hat r_{a}-r_{a}|\le\varepsilon_{R}\Bigr\}
\cap\Bigl\{\sup_{a\in\Acal}|\hat C_{a}-C_{a}|\le\varepsilon_{C}\Bigr\}
\label{eq:eventG}
\end{equation}
with probability at least $1-\alpha_{\mathrm{sel}}$. Screen actions by $\hat r_{a}+\varepsilon_{R}\le\delta$,
choose the minimum estimated-cost retained action, and certify only the frozen choice on independent data at
level $\alpha_{\mathrm{cert}}$, with $\alpha_{\mathrm{sel}}+\alpha_{\mathrm{cert}}\le\alpha$. Then
\begin{equation}
\Pr\bigl[r_{E_{t}}^{\hat a}\le\delta\ \ \text{or a valid fallback or abstention is invoked}\bigr]\ \ge\ 1-\alpha.
\label{eq:thm2}
\end{equation}
On the estimation event, a minimum-cost safe action with margin is retained and its cost is within
$2\varepsilon_{C}$ of optimal; a unique safe optimum with a larger cost gap is recovered exactly. This is a
sufficient protocol plus necessary obstruction statements, not an if-and-only-if characterization.
Section~\ref{sec:constructive} audits conditions (i)--(v) per target: (i), (ii), (iv), (v) are satisfiable
on the registered grid; (iii) fails for every one of the 48 targets, because every non-maximal action's risk
sits inside the margin band---a measured boundary, not a claim that certification is impossible.

\subsection{Certification scale and economic feasibility}

Classical uniform concentration gives the sufficient order
\begin{equation}
m=O\!\Bigl(\frac{\log(M/\alpha)}{\gamma^{2}}\Bigr).
\label{eq:mscale}
\end{equation}
With zero observed failures and per-action level $\alpha_l$, exact one-sided Clopper--Pearson certification
requires
\begin{equation}
m_{\min}=\Bigl\lceil\frac{\log\alpha_{l}}{\log(1-\delta)}\Bigr\rceil,
\label{eq:mmin}
\end{equation}
so 59 queries certify a single frozen action, but selection over $M=6$ actions under Bonferroni requires
94---and zero-failure certification passes with probability $(1-r)^m$ at true risk $r$, which is why the
maximum-budget action (risk $0.008$--$0.013$) certifies only 47--63\% of the time at $m=59$ in replay.
Query count, environment count, action count, grid resolution, and service workload control different
errors and are not interchangeable. For economics, let
\begin{equation}
D=G_{\mathrm{online}}-\Pr(R)\,\Delta C_{f\!-\!c}.
\label{eq:dgain}
\end{equation}
If the per-query net gain is nonpositive, no finite workload breaks even; otherwise the threshold is
\begin{equation}
N^{\star}=\frac{C_{\mathrm{build}}+C_{\mathrm{operator}}+C_{\mathrm{truth}}+C_{\mathrm{cert}}}{D}.
\label{eq:breakeven}
\end{equation}
This is an accounting identity, not an efficiency theorem. Mean gain cannot imply p95 or p99 improvement.

\begin{table}[t]
\centering
\caption{Theoretical results, assumptions, novelty boundary, and their empirical instantiation.}
\label{tab:theory}
\footnotesize
\setlength{\tabcolsep}{3.5pt}
\begin{tabular}{p{2.5cm}p{3.3cm}p{3.5cm}p{3.5cm}}
\toprule
Result & Core assumptions & Guarantee & Status of tools \\
\midrule
Transcript lower bound & conflicting regions; probe laws & four-outcome loss bound & classical testing; premise measured (\S\ref{sec:scale}) \\
Conditional recovery & feasibility, identifiability, & fixed-target familywise & new domain combination; margin \\
 & margin, evidence, fallback & safety + near-optimal cost & audited per target (\S\ref{sec:constructive}) \\
Certification scale & frozen $M$, margin $\gamma$ & sample order; zero-failure $m_{\min}$ & classical concentration \\
Economic gate & complete costs & finite break-even iff $D>0$ & accounting identity \\
\bottomrule
\end{tabular}
\end{table}

\section{ICBA: Certified Audit and Actionable Mitigation}
\label{sec:icba}

ICBA receives registered source and target builds, implementation-native actions, disjoint
design/selection/certification/evaluation roles, exact truth for authorized labeled roles, one
endpoint-aware failure event, implementation-internal costs, risk parameters, and a frozen candidate
family. It returns one of four scientifically distinct outcomes: certified reuse, target-specific
recalibration/certification, execution of an independently justified fallback, or abstention. It never
treats the maximum registered action as safe without evidence.

The procedure first validates role disjointness, replay, native action semantics, and cost fields. It then
constructs the categorical safe-budget response without endpoint imputation and measures directed
source-to-target transport. Any candidate selected using authorized target data is frozen before
independent certification, with familywise control across the registered actions. Finally, ICBA audits mean
cost, p95/p99 behavior, offline truth and control cost, fallback rate, and break-even separately. Missing
costs remain not estimable rather than being silently replaced by a proxy.

\begin{figure}[t]
\centering
\includegraphics[width=.8\linewidth]{fig2_workflow.png}
\caption{ICBA workflow. The procedure validates semantics before estimating transport, separates selection
from certification, and treats fallback and abstention as distinct terminal outcomes.}
\label{fig:workflow}
\end{figure}

A complementary mitigation removes, rather than learns, the build environment. The deterministic-rebuild
contract fixes the input permutation, construction seed where exposed, implementation and toolchain, and
requires single-threaded (or demonstrably deterministic-parallel) construction plus byte-level index replay
and search-output identity checks. Under these registered conditions the contract mechanically reuses the
same serialized environment; it is not an unseen-build guarantee and does not cover data refreshes, version
changes, or hardware changes. ICBA therefore implements a cost-ordered decision rule:
(1)~enforce a verified deterministic contract when feasible;
(2)~otherwise deploy robust source pooling when the operational history retains enough prior builds
(uncertified, risk below 1\% at 100K with 22 sources; scale-dependent);
(3)~otherwise collect target evidence and certify a frozen action, accepting that on grids like ours only
the maximum budget clears certification;
(4)~abstain. When near-zero risk is not required, canonical-order pinning under ordinary parallel construction is the
recommended \emph{default}: it halves response variation at no build-time cost ($0.80$--$0.93\times$),
leaving 30--34\% residual variation; the full single-threaded contract is reserved for settings that
require measured transport risk zero (Section~\ref{sec:contract}).

\section{Experimental Protocol}
\label{sec:protocol}

We evaluate SIFT-100K (128-dimensional L2 vectors) and Arxiv-Nomic-100K (768-dimensional normalized
embeddings under the registered L2/cosine-equivalent preprocessing) at Recall@10 with hit requirement
$h=10$, i.e., the returned top-10 must contain all ten ground-truth neighbors. hnswlib and the clean
Faiss-100K experiments contain 24 registered builds per dataset and all 552 non-diagonal directed pairs;
the Vamana-style stage contains 12 builds and 36 preregistered source-to-held-out-target directions. Each
cell uses 750 evaluation queries. Native action grids are reported per implementation (hnswlib
10/20/40/80/120/200; Faiss 16/32/64/128/256/512); raw \texttt{efSearch} and $l$ values are never equated
across systems. For Vamana, beam width is fixed to one, so its evidence is explicitly configuration
conditional.

\paragraph{SIFT-1M replication cell (preregistered).}
Eight builds (4 seeds $\times$ 2 insertion orders; $M{=}16$, efConstruction${}=100$, single-threaded,
mirroring the registered 100K configuration), a 500-query evaluation role sampled with seed 991 from the
ann-benchmarks test split, exact top-10 truth from the published ground truth, the registered 100K action
grid, and all metrics frozen in a manifest committed before any build or search. A query/base
content-overlap gate (zero required) passed before any risk number was computed; two additional
same-configuration builds verified byte-identity of the contract at 1M.

\paragraph{Data hygiene.}
Evaluation queries are unique and content-hash disjoint from the indexed base and historical analysis roles
in the clean Faiss-100K rerun; the Vamana-style stages receive the same audit in this revision: query/base
ID, raw-content, and normalized-content overlap, internal duplicates, and cross-stage leakage against the
hnswlib base are all zero on both datasets (historical-role raw-content overlap remains not estimable
because role vectors were not retained; the Vamana base is a separately registered snapshot, consistent
with never claiming cross-family numerical equality). Serialized-index replay, identifier mapping, and
native/tracer agreement are checked before inference. No future-replication truth or outcome is used; any
hash-only access status is retained in the artifact audit rather than promoted to a scientific claim.

Primary intervals use 5{,}000 query-bootstrap repetitions with seed 991 while retaining each sampled
query's full vector of directed-pair observations. The intervals are conditional on the registered builds,
not confidence intervals over unseen rebuilds. The primary estimand is the stage-registered endpoint-aware
incremental transport risk (Eq.~\ref{eq:delta}). For the four HNSW cells, the clean evidence supports a
common source-action convention; the Vamana stage used a different reference construction, and its frozen
records cannot reconstruct the fully harmonized HNSW estimand, so we report Vamana as supporting
stage-specific evidence rather than merging all six cells into one numerical range.

\section{Results}
\label{sec:results}

\subsection{Strictly comparable HNSW transport evidence}

The clean hnswlib and Faiss experiments agree on the main conclusion. Across the four data--implementation
cells, 75.87--91.07\% of evaluation queries exhibit finite-action safe-budget variation across registered
builds, whereas endpoint-state variation is only 0.53--2.67\%. Incremental source-to-target transport risk
ranges from 17.17\% to 23.60\% (Table~\ref{tab:crossfamily}). Every 95\% query-bootstrap interval excludes
the preregistered 2\% materiality threshold. Thus the primary signal is switching among finite native
actions, not a numerical artifact created by unresolved endpoints.

\begin{table}[t]
\centering
\caption{Registered cross-family evidence at Recall@10 with $h{=}10$. HNSW rows share the repaired common
estimand; the estimand crosswalk in Appendix~\ref{app:crossfamily} tabulates the coexisting registered
values.}
\label{tab:crossfamily}
\footnotesize
\setlength{\tabcolsep}{2.6pt}
\begin{tabular}{llccccc}
\toprule
Impl. & Dataset & B/p & $V_{\mathrm{fin}}$ & Risk [95\% CI] & $V_{\mathrm{end}}$ & Within-impl.\ cost \\
\midrule
hnswlib & SIFT-100K & 24/552 & 89.33\% & 21.55\% {\tiny[20.76, 22.33]} & 2.67\% & DistComp $+19.45\%$ \\
hnswlib & Arxiv-100K & 24/552 & 75.87\% & 17.17\% {\tiny[16.33, 18.06]} & 2.00\% & DistComp $+14.88\%$ \\
Faiss & SIFT-100K & 24/552 & 91.07\% & 23.60\% {\tiny[22.78, 24.39]} & 0.53\% & N/E (batch-cum.) \\
Faiss & Arxiv-100K & 24/552 & 75.87\% & 17.77\% {\tiny[16.86, 18.68]} & 1.87\% & N/E (batch-cum.) \\
\bottomrule
\end{tabular}
\end{table}

\begin{figure}[t]
\centering
\includegraphics[width=.9\linewidth]{fig3_crossfamily.png}
\caption{Cross-family evidence. Bars show safe-budget response variation; points show incremental transport
risk with 95\% query-bootstrap intervals. Vamana marks (daggered in Table~\ref{tab:vamana}) retain their
stage-specific estimand and are supporting rather than harmonized evidence.}
\label{fig:crossfamily}
\end{figure}

Robustness checks do not alter this conclusion. Removing the eight queries with the largest hnswlib risk
contribution leaves increments of 21.40\% and 16.90\%; deleting the top 1\% of Faiss risk contributors
leaves 23.43\% and 17.53\%; all 24 leave-one-build estimates remain positive on every cell. These are
robustness diagnostics conditional on the registered build family, not a population model for future
rebuilds.

\paragraph{Quality-threshold sensitivity.}
The $h{=}10$ event (all ten neighbors returned) is the strictest registered target, so we recomputed the
family transport risk under relaxed hit requirements $h\in\{8,9\}$ for every cell
(Table~\ref{tab:hsens}). Risk decreases monotonically as the requirement relaxes, but the minimum over
all cells at $h{=}8$ is 5.96\%---still three times the 2\% materiality gate. The phenomenon is not an
artifact of near-exact retrieval.

\begin{table}[t]
\centering
\caption{Transport risk under relaxed quality thresholds (family estimator of
Section~\ref{sec:results}; hnswlib-100K rows are registered values, others computed from the same
frozen records).}
\label{tab:hsens}
\footnotesize
\setlength{\tabcolsep}{4pt}
\begin{tabular}{lccc}
\toprule
Cell & $h{=}8$ & $h{=}9$ & $h{=}10$ \\
\midrule
hnswlib SIFT-100K & 15.60\% & 19.10\% & 22.03\% \\
hnswlib Arxiv-100K & 9.36\% & 13.39\% & 17.60\% \\
Faiss SIFT-100K (clean) & 10.56\% & 16.81\% & 23.67\% \\
Faiss Arxiv-100K (clean) & 5.96\% & 10.32\% & 17.93\% \\
SIFT-1M (preregistered) & 19.77\% & 22.55\% & 27.95\% \\
\bottomrule
\end{tabular}
\end{table}

\subsection{Vamana-style supporting boundary}

\begin{table}[t]
\centering
\caption{Vamana-style boundary panel (dagger: stage-specific estimand, beam width one, asymmetric
directions; not harmonized with the HNSW estimand and not merged into its range).}
\label{tab:vamana}
\footnotesize
\setlength{\tabcolsep}{2.6pt}
\begin{tabular}{llccccc}
\toprule
Impl. & Dataset & B/p & Cat.\ var. & Risk [95\% CI] & Cens. & Cost (descr.) \\
\midrule
DiskANN3/Vam.$^\dagger$ & SIFT-100K & 12/36 & 71.07\% & 16.86\% {\tiny[15.65, 18.10]} & 0.27\% & $\sim$1.00\% {\tiny[$-$0.94, 3.04]} \\
DiskANN3/Vam.$^\dagger$ & Arxiv-100K & 12/36 & 48.27\% & 11.37\% {\tiny[10.20, 12.60]} & 0.80\% & $\sim$1.11\% {\tiny[$-$1.05, 3.22]} \\
\bottomrule
\end{tabular}
\end{table}

The Vamana-style stage reproduces the risk direction on both datasets under its original stage-specific
endpoint estimand. It uses 12 builds, 36 asymmetric directions, beam width one, and a reference construction
that cannot be harmonized exactly from the frozen event table. It therefore broadens the
construction-family scope without supporting numerical equality to HNSW. Cost is even less universal:
hnswlib shows resolved within-implementation DistComp penalties of 19.45\% and 14.88\%, whereas the
Vamana-style estimates are near 1\% with intervals crossing zero, and clean Faiss-100K distance-computation
cost is not estimable because the exposed counter is batch cumulative.

\subsection{Deterministic rebuilding removes the registered environment}
\label{sec:contract}

We fix the input order, seed, thread count, implementation, and toolchain, and require byte-identical
serialization and search replay. Three repeated builds per dataset satisfy all identity checks and have
measured transport risk zero. Relative to the original multithreaded construction, the contract preserves
Recall@10 within 0.0004, leaves p95/p99 safe budgets unchanged, changes mean safe budgets by ratios 0.976
(SIFT) and 0.999 (Arxiv), and increases build time by $3.52\times$ and $1.71\times$. This is an actionable
environment-elimination contract---a hnswlib deterministic construction control---not a recovery algorithm
for changed data or software.

\begin{table}[t]
\centering
\caption{Contract outcomes and primitive profiling economics (8-thread resident measurements; primitive
operations only, not a deployment ledger).}
\label{tab:contract}
\footnotesize
\setlength{\tabcolsep}{3.5pt}
\begin{tabular}{p{3.7cm}llp{4.2cm}}
\toprule
Result & SIFT-100K & Arxiv-100K & Interpretation \\
\midrule
Deterministic transport risk & 0 & 0 & registered same-data/toolchain contract \\
Recall difference vs.\ original & $+0.00027$ & $-0.00040$ & within the registered $\pm0.001$ tolerance \\
Mean safe-budget ratio & 0.976 & 0.999 & descriptive common-feasible summary \\
p95 / p99 budget ratio & 1.00 / 1.00 & 1.00 / 1.00 & no measured tail change \\
Build-time ratio vs.\ original & $3.52\times$ & $1.71\times$ & determinism has offline overhead \\
59-query primitive profiling & H 0.117\,s; F 0.017\,s & H 0.994\,s; F 0.248\,s & cheap, operator dependent \\
\bottomrule
\end{tabular}
\end{table}

\paragraph{Chain ablation (descriptive).}
A four-tier chain---D0C (random order, 8 threads, random seeds), D1 (canonical order pinned), D2 (seed
pinned), D3 (single-threaded)---attributes the contract's effect: pinning the insertion order to a
canonical order halves minimum-safe-action variation (68.3 to 34.1\% on SIFT; 63.2 to 30.8\% on Arxiv) at
build-time ratios $0.80\times$ and $0.93\times$, i.e., free. Pinning the seed under 8-thread construction is
nearly inert ($-1.3$ and $-4.3$ points), and the six D2 builds with identical order and seed produce six
distinct serialized indexes: parallel scheduling is a residual nondeterminism source that the exposed seed
does not control. Single-threaded construction is therefore the necessary and sufficient registered clause
for byte identity (3/3 identical, variation 0) and carries the entire overhead ($4.73\times$ and
$1.86\times$ relative to the pinned 8-thread tier; $3.52\times$ and $1.71\times$ relative to D0C). The
ablation is a chain, not a full factorial; effects are descriptive.

\begin{figure}[t]
\centering
\includegraphics[width=.95\linewidth]{fig6_ablation.png}
\caption{Deterministic-contract ablation (descriptive chain). Left bars: minimum-safe-action variation per
tier; right line: median build time. Order pinning halves variation for free; single-threading removes all
variation and carries the build overhead.}
\label{fig:ablation}
\end{figure}

\subsection{Constructive baselines and certification}
\label{sec:constructive}

We evaluate six policies on the hnswlib cells by retrospective replay over the frozen per-query records
(roles 375/94/281 disjoint, seed 991; bootstrap as registered). Table~\ref{tab:constructive} reports
target-truth usage, risk under the four-outcome semantics, DistComp versus the per-query oracle, and
certification outcomes.

\begin{table}[t]
\centering
\caption{Constructive decision table (hnswlib, registered grid, retrospective replay). M2 dominates on risk
at $1.34$--$1.45\times$ oracle cost but requires 22 retained source builds and is uncertified; certified
routes collapse to the maximum budget.}
\label{tab:constructive}
\scriptsize
\setlength{\tabcolsep}{2.8pt}
\begin{tabular}{p{3.2cm}p{1.0cm}p{3.0cm}p{1.4cm}p{4.0cm}}
\toprule
Method & Truth & Risk (SIFT / Arxiv) & DistC. & Certification \\
\midrule
M1 naive transport ($k{=}1$) & 0 & 21.89\% / 18.06\% & 2.6/2.7$\times$ & none \\
M2 max-over-source ($k{=}22$) & 0 & 0.18\% / 0.09\% (abst.\ 2.4/2.7\%) & 1.45/1.34$\times$ & uncert.; LOBO $\le$2.1\% \\
M3 frozen max-action CP ($M{=}1$, $m{=}59$) & 59 q & 0.87\% / 1.36\% & 4.16/5.28$\times$ & cert.\ 71\%/38\% of draws \\
M4a point-screen sel.-then-cert. & 469 q & --- & --- & 0/48 (zero-fail.\ violated) \\
M4b Theorem-2 screen, then certify & 469 q & --- & --- & 0/48; margin too thin \\
M5 always maximum action & 0 & 0.80\% / 1.26\% & 4.09/5.14$\times$ & not auto.\ safe \\
M6 per-query oracle & full & endpoint mass only & 1.00$\times$ & non-deployable ref. \\
\bottomrule
\end{tabular}
\end{table}

\begin{figure}[t]
\centering
\includegraphics[width=.95\linewidth]{fig4_pooling.png}
\caption{Robust source pooling: transport risk versus number of pooled source builds (mean and p95 over 50
draws; log scale). $k{=}1$ reproduces the registered naive-transport risks; $k\ge10$ falls below 2.4\% at
100K. At SIFT-1M the decay is slower ($k{=}7$ still leaves 10.4\%): the pooling requirement grows with
scale.}
\label{fig:pooling}
\end{figure}

\begin{figure}[t]
\centering
\includegraphics[width=.75\linewidth]{fig5_plane.png}
\caption{Constructive decision plane (hnswlib). Risk versus DistComp against the per-query oracle;
horizontal gates at 2\% (materiality) and 5\% (risk tolerance). M2 occupies the low-risk/low-cost corner
without certification; certified M3 coincides with M5.}
\label{fig:plane}
\end{figure}

\paragraph{Theorem-2 verdict.}
Conditions (i), (ii), (iv), (v) are satisfiable on the registered grid; (iii) fails for all 48 targets:
every non-maximal action's risk lies inside the margin band, so Theorem-2 screening passes (it selects
$\mathrm{ef}{=}120$ on eight targets and $\mathrm{ef}{=}200$ on sixteen per dataset) but zero-failure
certification then fails on every target. The certified route collapses to
maximum-budget-or-abstain, and even the maximum action certifies only with probability 0.47--0.63 at
$m=59$ because its residual risk (0.008--0.013, concentrated in always-infeasible queries) violates the
zero-failure requirement in roughly half the draws. A $\gamma$-sensitivity audit sharpens the verdict: at
$\gamma{=}0.005$ a non-maximal action with selection-block risk $\le\delta-2\gamma$ \emph{exists} on
16/24 (SIFT) and 21/24 (Arxiv) targets (16/24 and 8/24 at $\gamma{=}0.01$; none at $\gamma{=}0.02$), so
the margin condition is not absolutely absent---what fails everywhere is the \emph{certification step}:
margin-band actions carry 3--4\% risk and pass zero-failure certification at $m{=}94$ with probability
$(1-r)^{94}\approx2$--$6\%$. Certification power, not margin existence, is the binding constraint at
registered sample sizes. This is a measured boundary of the recovery theorem on this grid---conditional,
not an impossibility claim. Pooling is the constructive winner when operational
history retains enough prior builds; its uncertainty remains a registered-family bootstrap, not an
unseen-build guarantee. Non-max aggregation does not help: at $k{=}22$ the $0.9$-quantile of source
actions is nearly identical to the max (1.07\% vs.\ 0.96\% risk at 45.5\% vs.\ 48.4\% conservatism on
SIFT), while median pooling cuts conservatism to 16--22\% but leaves 13--21\% risk---the response
diameter is driven by outlier sources, so there is no cheaper order-statistic operating point.

\subsection{Scale, learned-predictor transfer, and distinguishability}
\label{sec:scale}

\begin{table}[t]
\centering
\caption{Preregistered P6 probes: scale replication, pooling boundary, predictor transfer, and transcript
distinguishability.}
\label{tab:p6}
\footnotesize
\setlength{\tabcolsep}{3pt}
\begin{tabular}{p{3.2cm}p{6.3cm}p{4.2cm}}
\toprule
Probe & Result & Boundary closed \\
\midrule
SIFT-1M cell (prereg.; 8 builds, 500 q) & incr.\ risk 21.87\% [21.52, 22.21]; $V_{\mathrm{fin}}$ 98.4\% vs $V_{\mathrm{end}}$ 13.2\%; contract byte-identical; profiling 0.067\,s & phenomenon replicates at $10\times$ scale (100K: 21.55\%) \\
Pooling at 1M & $k{=}7$ leaves 10.4\% risk (100K w/ 22 sources: $<1\%$); quantile pooling: $q{=}0.9 \approx$ max, median unusable & pooling requirement scale-dependent \\
Learned predictor (HistGB, ef{=}10 transcript) & transfer 23.9\%/21.2\% vs naive 21.1\%/17.3\%; in-target CV 24.1\%/22.2\%; $R^2=0.015$ & layer-3 non-portability; $B(q)$ build-structure dominated \\
Distinguishability (552 pairs, $k\in\{1,3,6\}$) & hits: acc.\ 0.49--0.51, TV $\le0.25$; runtime features: acc.\ 0.56--0.58 (max 0.72) & Thm.~1 premise holds for hit transcripts; runtime probes partially separate \\
\bottomrule
\end{tabular}
\end{table}

The predictor probe deserves emphasis: a regressor trained on 23 source builds predicts the target's
minimum safe action no better than chance ($R^2=0.015$), and even in-target cross-validation leaves
22--24\% risk---the deployment-time transcript of the cheapest action carries almost no information about
which build-specific neighborhood structure makes a larger budget necessary. The three-layer boundary is
now complete: response labels are strongly non-portable (layer 1), global policies have a constructive
bounded answer (layer 2, Section~\ref{sec:constructive}), and learned per-query predictors do not transfer
on these features (layer 3). We did not test richer probe features or model classes; the claim is
feature-class-conditional.

\subsection{Economics ledger and grid sensitivity}

\begin{table}[t]
\centering
\caption{Six-cell economics matrix (selected rows; the full 36-cell matrix with per-cell NOT\_ESTIMABLE
reasons is in the artifact). N/E = not estimable from frozen evidence; never imputed.}
\label{tab:economics}
\footnotesize
\setlength{\tabcolsep}{3.4pt}
\begin{tabular}{lccccc}
\toprule
Component & hnS & hnA & FS & FA & Vam \\
\midrule
Build time (D3 median) & 12.9\,s & 46.6\,s & N/E & N/E & N/E \\
Profiling, resident 59q (8t) & 0.117\,s & 0.994\,s & 0.017\,s & 0.248\,s & N/E \\
Break-even, search-only & NO\_FINITE & 11{,}832\,q & N/E & N/E & N/E \\
Break-even, wall-clock & NO\_FINITE & 104{,}383\,q & N/E & N/E & N/E \\
Candidate replay / control & N/E & N/E & N/E & N/E & N/E \\
\bottomrule
\end{tabular}
\end{table}

Transport risk is grid-indexed: on registered subgrids (drop-min, drop-max, coarse 3-action) mean family
risk stays between 10.3\% and 23.8\%---always above the 2\% materiality gate---while coarser grids
mechanically lower measured risk because the minimum safe action can only grow. We therefore claim
phenomenon-level grid-robustness, not resolution-invariant risk levels, and keep native grids separate
across implementations.

\section{Theory-Guided Falsification Ladder}
\label{sec:ladder}

The transcript-limited lower bound identifies insufficient build information as one possible barrier, while
Theorem~2 states what a fixed-target escape requires. We therefore organize method development as a
falsification ladder rather than a performance tournament. A route advances only if its semantics,
candidate attainability, observable information or structural bridge, independent certification, tail
behavior, and deployment value survive in that order.

\begin{table}[t]
\centering
\caption{Four recovery families and the first prerequisite falsified by sealed evidence; the 16-route
casebook is in Appendix~\ref{app:ladder}.}
\label{tab:ladder}
\footnotesize
\setlength{\tabcolsep}{3pt}
\begin{tabular}{p{2.3cm}p{4.2cm}p{4.2cm}p{2.7cm}}
\toprule
Recovery family & Strongest positive evidence & Decisive limitation & Inference \\
\midrule
Reuse / recalibration & fixed-target safety in selected designs; & no stable history increment; & fixed-target safety only \\
 & probe quantifies the negative side ($R^2{=}0.015$) & disjoint eval.\ removes gain & \\
Build selection / stable constr. & safe-selection theory valid; isolated gains & no CI-supported joint action (72 frozen) & generator closed \\
Multi-lane / portal rescue & complementarity replicates after ID repair & no truth-free trigger; 79.5--93.4\% compression & mechanism only \\
Certified auditor / fallback & 36/36 accepted decisions safe & collapse to max budget; regret 470.3 / 684.5 & diagnostic, not optimizer \\
\bottomrule
\end{tabular}
\end{table}

\begin{figure}[t]
\centering
\includegraphics[width=.95\linewidth]{fig7_ladder.png}
\caption{Theory-guided falsification ladder. Each recovery family is pruned at its first unsupported
prerequisite. Partial mechanism evidence is retained without being promoted to a deployable algorithm.}
\label{fig:ladder}
\end{figure}

Target-only residual recalibration initially reduced mean distance computations by 29.40\% (SIFT) and 26.17\%
(Arxiv) in a paired fixed-target design; history-assisted variants added only 0.90\% [$-8.73$, 11.25] and
$-0.13$\% [$-9.24$, 10.26]; under strict disjoint-query cross-fitting the target-only effects became
$-7.35$\% and $-12.49$\%. Build selection failed earlier: among 72 frozen build-budget actions, SIFT had no
point-estimate-or-interval supported jointly feasible action and Arxiv had one point-estimate action with no
interval support. One protected-edge repair preserved graph invariants yet reduced Recall@10 by 0.002 and
worsened both endpoint families. Fixed portal sets showed genuine rescue complementarity (1{,}064 and 727
additional hits; 426 and 348 threshold rescues) but cost $2.5$--$2.8\times$ a matched native lane and
remained 3.3--4.8 recall points lower; eight preregistered truth-free triggers showed no positive validation
advantage, and a shared-frontier bound implies 79.5--93.4\% auxiliary-cost compression before viability
(provenance for these aggregates is itemized in the artifact; three of them are recorded as
origin-unlocated and are being re-derived from row-level records). Finally, the conservative auditor
demonstrated safety without value across 60 retrospective decisions (4 direct reuses, 32 certified
fallbacks, 24 abstentions; all 36 accepted decisions safe, every accepted action at the maximum registered
budget).

\section{Limitations and Scope}
\label{sec:limitations}

Our primary evidence covers two 100K-scale datasets and four comparable HNSW cells, plus one preregistered
1M hnswlib replication; the 1M cell is single-implementation and eight builds. The Vamana extension uses a
different estimand, asymmetric pairs, and beam width one. We do not model future builds, data refreshes, or
open-world safety. Scale, pruning, I/O, and hardware may change effect sizes; clean Faiss
distance-computation cost and Vamana wall-clock/I/O are not estimable, precluding a universal cost claim.

The lower bound is conditional and probe-class-limited; our instantiation covers the registered hit-count
transcript class, and other probe classes (latency, visited counts) may separate builds better. The recovery
theorem requires a useful margin-separated action, which the falsification ladder and the per-target audit
both reject on this grid; richer candidate families might restore it. The 2\% gate is a preregistered
scientific threshold, not a deployment utility function.

The deterministic contract eliminates registered same-data variation but costs $1.71$--$3.52\times$ more
build time and does not cover changed data, software, or operators; the ablation attributing its effect is a
descriptive chain, not a factorial. At both measured scales primitive profiling is cheap; the decision rule,
not profiling cost, carries the practical weight. The constructive table is a retrospective replay on
registered builds: pooling requires retaining prior builds and their per-query records, and its risk
estimate is an uncertified registered-family bootstrap whose source requirement grows with scale ($k{=}7$
leaves 10.4\% at 1M). The predictor probe is one model class on one feature family. Any larger-scale
advantage of reuse remains unverified until truth, control, latency, and I/O costs are measured end to end.

The historical baseline remains conditionally reproducible (111/123 matching entries). Later clean reruns do
not erase this limitation. The artifact records hash-only role access, Vamana reference semantics, and the
descriptive---not causal---status of construction-factor comparisons.

\section{Conclusion}
\label{sec:conclusion}

Safe search budgets are properties of a realized Graph-ANNS build, not only of a dataset and nominal
recipe. Across two datasets, two HNSW implementations, and a preregistered 1M replication, rebuilds cause
large finite-action variation and 17--24\% incremental transport risk; a Vamana-style study supports the
direction under a distinct estimand. The theory separates transcript-limited non-portability---whose
near-indistinguishability premise we measure directly---from fixed-target recovery, whose margin condition
we find systematically absent on the registered grid. Practically, the evidence orders itself into a
decision rule: verify a deterministic single-threaded rebuild contract when feasible; otherwise pool many
retained source builds, accepting an uncertified but measured risk below 1\% at 100K and a scale-dependent
source requirement; otherwise certify a frozen action with declared margin, expecting only the maximum
budget to clear on grids like ours; otherwise abstain. Learned per-query predictors do not close this gap on
deployment-time transcript features. The boundary between what transports and what must be re-established
per build is now measured at every layer we probe.

\begin{thebibliography}{99}
\itemsep0pt

\bibitem[Angelopoulos et~al.(2025)]{angelopoulos2025ltt}
Angelopoulos, A.~N., Bates, S., Candes, E.~J., Jordan, M.~I., and Lei, L.
Learn then test: Calibrating predictive algorithms to achieve risk control.
\emph{Annals of Applied Statistics}, 2025.

\bibitem[Angelopoulos et~al.(2024)]{angelopoulos2024crc}
Angelopoulos, A.~N., Bates, S., Fisch, A., Lei, L., and Schuster, T.
Conformal risk control.
\emph{International Conference on Learning Representations}, 2024.

\bibitem[Bae et~al.(2026)]{bae2026qbat}
Bae, J., Ham, T.~J., Li, A., Chockchowwat, S., and Papakonstantinou, Y.
QBAT: Model-based query budget autotuner for clustering-based approximate nearest neighbor search.
\emph{Proceedings of the VLDB Endowment}, 2026.

\bibitem[Bates et~al.(2021)]{bates2021rcps}
Bates, S., Angelopoulos, A.~N., Lei, L., Malik, J., and Jordan, M.~I.
Distribution-free, risk-controlling prediction sets.
\emph{Journal of the ACM}, 68(6), 2021.

\bibitem[Ben-David et~al.(2010)]{bendavid2010theory}
Ben-David, S., Blitzer, J., Crammer, K., Kulesza, A., Pereira, F., and Vaughan, J.~W.
A theory of learning from different domains.
\emph{Machine Learning}, 79:151--175, 2010.

\bibitem[Chatzakis et~al.(2025)]{chatzakis2025darth}
Chatzakis, M., Papakonstantinou, Y., and Palpanas, T.
DARTH: Declarative recall through early termination for approximate nearest neighbor search.
\emph{Proceedings of the ACM on Management of Data}, 3(4), 2025.

\bibitem[Duchi et~al.(2025)]{duchi2025multienv}
Duchi, J.~C., Gupta, S., Jiang, K., and Sur, P.
Predictive inference in multi-environment scenarios.
\emph{Statistical Science}, 40(3):392--416, 2025.

\bibitem[Elliott and Clark(2024)]{elliott2024order}
Elliott, O.~P. and Clark, J.
The impacts of data, ordering, and intrinsic dimensionality on recall in hierarchical navigable small
worlds. \emph{arXiv:2405.17813}, 2024.

\bibitem[El-Yaniv and Wiener(2010)]{elyaniv2010selective}
El-Yaniv, R. and Wiener, Y.
On the foundations of noise-free selective classification.
\emph{Journal of Machine Learning Research}, 11:1605--1641, 2010.

\bibitem[Fu et~al.(2019)]{fu2019nsg}
Fu, C., Xiang, C., Wang, C., and Cai, D.
Fast approximate nearest neighbor search with the navigating spreading-out graph.
\emph{Proceedings of the VLDB Endowment}, 12(5):461--474, 2019.

\bibitem[Garivier and Kaufmann(2016)]{garivier2016bai}
Garivier, A. and Kaufmann, E.
Optimal best arm identification with fixed confidence.
\emph{Proceedings of the 29th Conference on Learning Theory}, pages 998--1027, 2016.

\bibitem[Horchidan et~al.(2025)]{horchidan2025conann}
Horchidan, S., Zeiher, F., Bostrom, H., and Carbone, P.
ConANN: Conformal approximate nearest neighbor search.
\emph{Proceedings of the VLDB Endowment}, 19(1):29--42, 2025.

\bibitem[Johansson et~al.(2019)]{johansson2019support}
Johansson, F.~D., Sontag, D., and Ranganath, R.
Support and invertibility in domain-invariant representations.
\emph{Proceedings of the 22nd International Conference on Artificial Intelligence and Statistics}, pages
527--536, 2019.

\bibitem[Johnson et~al.(2021)]{johnson2021faiss}
Johnson, J., Douze, M., and Jegou, H.
Billion-scale similarity search with GPUs.
\emph{IEEE Transactions on Big Data}, 7(3):535--547, 2021.

\bibitem[Li et~al.(2020)]{li2020learned}
Li, C., Zhang, M., Andersen, D.~G., and He, Y.
Improving approximate nearest neighbor search through learned adaptive early termination.
\emph{Proceedings of the 2020 ACM SIGMOD International Conference on Management of Data}, pages
2539--2554, 2020.

\bibitem[Malkov and Yashunin(2020)]{malkov2020hnsw}
Malkov, Y.~A. and Yashunin, D.~A.
Efficient and robust approximate nearest neighbor search using hierarchical navigable small world graphs.
\emph{IEEE Transactions on Pattern Analysis and Machine Intelligence}, 42(4):824--836, 2020.

\bibitem[Prokhorenkova and Shekhovtsov(2020)]{prokhorenkova2020graph}
Prokhorenkova, L. and Shekhovtsov, A.
Graph-based nearest neighbor search: From practice to theory.
\emph{Proceedings of the 37th International Conference on Machine Learning}, pages 7803--7813, 2020.

\bibitem[Subramanya et~al.(2019)]{subramanya2019diskann}
Subramanya, S.~J., Devvrit, Simhadri, H.~V., Krishnaswamy, R., and Kadekodi, R.
DiskANN: Fast accurate billion-point nearest neighbor search on a single node.
\emph{Advances in Neural Information Processing Systems}, 32, 2019.

\bibitem[Wang et~al.(2026a)]{wang2026annie}
Wang, Z., Chatzakis, M., Wang, Q., Palpanas, T., Wang, P., and Wang, W.
ANNiE: A learned query cost estimator for graph-based approximate nearest neighbor search.
\emph{Proceedings of the VLDB Endowment}, 19(11):3820--3833, 2026.

\bibitem[Wang et~al.(2022)]{wang2022safebai}
Wang, Z., Wagenmaker, A., and Jamieson, K.
Best arm identification with safety constraints.
\emph{Proceedings of the 25th International Conference on Artificial Intelligence and Statistics}, 2022.

\bibitem[Xu et~al.(2022)]{xu2022maintenance}
Xu, Z., Zhao, W., Tan, S., Zhou, Z., and Li, P.
Proximity graph maintenance for fast online nearest-neighbor search.
\emph{arXiv:2206.10839}, 2022.

\bibitem[Yu(1997)]{yu1997assouad}
Yu, B.
Assouad, Fano, and Le Cam.
In \emph{Festschrift for Lucien Le Cam}, pages 423--435. Springer, 1997.

\bibitem[Zhang and Miller(2026)]{zhang2026exploration}
Zhang, C. and Miller, R.~J.
Distribution-aware exploration for adaptive HNSW search.
\emph{Proceedings of the ACM on Management of Data}, 4(1), 2026.

\end{thebibliography}

\appendix

\section{Complete Notation and Semantic Contracts}
\label{app:notation}

\begin{table}[h]
\centering
\caption{Core notation and scope guardrails.}
\small
\begin{tabular}{lll}
\toprule
Symbol & Meaning & Scope guardrail \\
\midrule
$\Ecal$ & preregistered build-environment set & no population over unseen builds \\
$E$ & one serialized build environment & includes replay semantics \\
$\Acal_E$ & finite native action set & implementation specific \\
$Z_E(q,a)$ & endpoint-aware failure indicator & shared across roles \\
$B_E(q)$ & minimum observed safe action & registered grid only \\
$\BOT$ & unresolved registered endpoint & never numerically imputed \\
$C_E(q,a)$ & distance computations & within one implementation \\
$\Delta_{s\to t}$ & incremental transport risk (Eq.~\ref{eq:delta}) & registered-family conditional \\
$V_{\mathrm{fin}},V_{\mathrm{end}}$ & variation families (Eq.~\ref{eq:var}) & reported separately \\
$\delta,\alpha,\gamma,M$ & 0.05; 0.05; margin; frozen family size & preregistered \\
\bottomrule
\end{tabular}
\end{table}

For HNSW, raw fixed \texttt{efSearch} is treated as a finite action rather than an assumed monotone
stopping sequence. For Vamana-style search, $l$ and beam width are distinct; beam width is fixed to one in
the registered experiment. Native budget and cost counters remain implementation-specific.

\section{Proof of Theorem 1}
\label{app:proof1}

Define an induced binary decision $I$: return 0 when the chosen action lies in $D_0$, return 1 when it lies
in $D_1$, and assign actions outside both regions arbitrarily. Since $D_0\cap D_1=\varnothing$, an action
outside $D_i$ incurs loss at least $\Delta$. Therefore
\begin{equation}
\E_{0}[l_{0}]+\E_{1}[l_{1}]\ \ge\ \Delta\bigl\{P_{0}[I{=}1]+P_{1}[I{=}0]\bigr\}.
\label{eq:proof1}
\end{equation}
The minimum sum of binary testing errors equals $1-\TV$; the maximum expected loss is at least half the sum,
proving Theorem~1. Bretagnolle--Huber lower-bounds the error sum and Pinsker yields the KL consequence.
For adaptive observations the chain rule gives
\begin{equation}
K_{m}=\sum_{t=1}^{m}\E_{0}\!\left[\KL\bigl(P_{0}^{Y_{t}\mid H_{t-1},A_{t}}\,\bigl\|\,P_{1}^{Y_{t}\mid H_{t-1},A_{t}}\bigr)\right]\le m\,I^{\star},
\label{eq:chain}
\end{equation}
and data processing through the decision rule yields the stated testing consequence. The testing
inequalities are classical; only the Graph-ANNS deployment-loss reduction is domain specific. The
instantiation in Section~\ref{sec:scale} supplies the measured TV bound the premise requires.

\section{Proof of Theorem 2}
\label{app:proof2}

On the estimation event every retained action satisfies the screened bound. If a minimum-cost safe action
$a_k$ has margin $r_{a_k}\le\delta-2\varepsilon_R$, then
\begin{equation}
\hat r_{a_{k}}+\varepsilon_{R}\ \le\ r_{a_{k}}+2\varepsilon_{R}\ \le\ \delta,
\label{eq:retain}
\end{equation}
so it is retained; since the deployed action $\hat a$ minimizes estimated cost among retained actions,
\begin{equation}
\hat C_{\hat a}\ \le\ \hat C_{a_{k}}+\varepsilon_{C}\ \le\ C_{a_{k}}+2\varepsilon_{C}.
\label{eq:cost}
\end{equation}
If every other safe action costs strictly more than the optimum plus $2\varepsilon_C$, these inequalities
force exact selection on the event. Selection and certification roles are independent and the final action
is frozen; a valid upper certificate plus a union bound gives the $1-\alpha$ guarantee. On rejection the
procedure executes only an independently valid fallback; without one it abstains. The endpoint obstruction
(every in-family policy fails on the unresolved or infeasible subset), the identifiability obstruction
(Theorem~1), the no-margin obstruction (risks $\delta-\epsilon$ and $\delta+\epsilon$ become
indistinguishable as $\epsilon\downarrow0$), adaptive selection bias, and rejection without a safe fallback
are necessary obstructions, not an if-and-only-if theorem.

\section{Corollaries and Restricted Propositions}
\label{app:corollaries}

\paragraph{Reverse-pair symmetry.} If every unordered pair appears in both directions with equal weight,
each unequal pair contributes one under-budget and one over-budget direction, hence
\begin{equation}
P_{\mathrm{under}}=P_{\mathrm{over}}=\tfrac{1}{2}\,P\bigl[B_{s}\neq B_{t}\bigr],
\label{eq:sym}
\end{equation}
a combinatorial identity; asymmetric source-to-held-out-target designs are still needed for directional
portability claims.

\paragraph{Endpoint decomposition.} With $\eta_E$ the mass on which no registered action succeeds,
\begin{equation}
r_{E}^{\pi}=\eta_{E}+(1-\eta_{E})\,r_{E}^{\pi,\mathrm{feas}},
\label{eq:endpoint}
\end{equation}
so $\eta_E>\delta$ precludes in-family certification.

\paragraph{Failure complementarity.} With conditional rescue rate $\rho_S$,
\begin{equation}
r_{S}=r_{0}\,(1-\rho_{S}),
\label{eq:rescue}
\end{equation}
and marginal auxiliary accuracy cannot replace the conditional rescue rate; shared queues, visited sets,
truncation, or altered tie rules can invalidate the mechanical preservation condition.

\paragraph{Structural bridges.} The only generally valid bound retained is
\begin{equation}
\bigl|r_{G',a}-r_{G,a}\bigr|\ \le\ \Pr_{q}\{Z_{G}(q,a)\neq Z_{G'}(q,a)\},
\label{eq:bridge}
\end{equation}
and expansion-prefix arguments require fixed search semantics, preserved critical edges, priority margins,
and bounded frontier intruders; they do not imply a theorem for raw \texttt{efSearch} or $l$ without an
empirical implementation bridge.

\section{Complete Experimental Protocol}
\label{app:protocol}

The registered unit is a serialized build, not a query row. The strict HNSW evidence uses 24 hnswlib and 24
clean Faiss-100K builds per dataset and all 552 non-diagonal directed pairs; the Vamana-style analysis uses
12 builds and 36 preregistered directions; the 1M cell uses 8 builds, 56 directed pairs, and 500
preregistered evaluation queries. Every 100K cell contains 750 evaluation queries at Recall@10 with $h=10$.
The clean Faiss query roles have zero ID, raw-content, and normalized-content overlap with the indexed base
and historical roles; the Vamana stages carry the same audit at zero overlap with the caveats stated in
Section~\ref{sec:protocol}. Native action grids and counters are never numerically equated across
implementations.

The query bootstrap resamples query identities while retaining the complete vector of registered
directed-pair observations; it quantifies query-distribution uncertainty conditional on the finite build
family. Leave-one-build-out and leave-one-seed/order analyses test domination but do not define an
unseen-build population. Top-contributor deletion removes 1\% of evaluation queries under a frozen
contribution rule (eight queries under the registered hnswlib ordering). The confidence intervals support
registered-family conclusions only. An unresolved endpoint remains $\BOT$ in all categorical analyses.
Mean, p95, p99, build cost, truth cost, control cost, and fallback cost are separate fields; missing items
remain not estimable.

\section{Cross-Family Details and Estimand Crosswalk}
\label{app:crossfamily}

In hnswlib, same-seed cross-order inconsistency is 57.43\% (SIFT) and 45.23\% (Arxiv), substantially larger
than same-order cross-seed inconsistency (13.25\% and 12.62\%); the comparison is descriptive because
factors are not independently randomized across implementations. Faiss does not expose a reliable
construction seed under the registered interface, so fixed input permutations define its 24 build
environments; the clean 100K rerun removes query/base self-matches. DiskANN3/Vamana-style preflight
establishes fixed-configuration replay and nonzero differences across registered permutations, but the
frozen Vamana event table cannot reconstruct the harmonized HNSW source-action estimand; Vamana results
remain stage specific. Finite-action and endpoint variation are reported separately (75.87--91.07\% versus
0.53--2.67\%); unresolved mass is 0.072/0.167\% (Faiss SIFT/Arxiv) and 0.767/1.067\% (hnswlib SIFT/Arxiv).

\paragraph{Estimand crosswalk (hnswlib).}
The E4 primary oracle-transport estimand gives 21.57\% [20.74, 22.40] (SIFT) and 17.12\% [16.20, 18.06]
(Arxiv) under crossed seed-query bootstrap; the original cross-family stage gives 21.57\% [20.77, 22.35]
and 17.12\% [16.24, 18.00]; the repaired $h{=}10$ hit-count common estimand---adopted in
Table~\ref{tab:crossfamily} because the clean Faiss evidence exists only under hit-count semantics---gives
21.55\% [20.76, 22.33] and 17.17\% [16.33, 18.06]. All three coexist as registered values; the paper adopts
the repaired common estimand for the four HNSW cells and does not modify any frozen number.

\section{Falsification-Ladder Casebook}
\label{app:ladder}

\begin{table}[h]
\centering
\caption{Sixteen route-level tests summarized by family and first failed condition. Three historical
aggregates (the recalibration percentages, the protected-edge repair cost pair, and the exact matched-lane
ratios) lack stored provenance and are flagged for re-derivation in the artifact.}
\small
\begin{tabular}{lll}
\toprule
Family & Route & First failed prerequisite \\
\midrule
Reuse/recal. & source-only global reuse & nontrivial portability \\
Reuse/recal. & target-only residual calibration & disjoint-query replication \\
Reuse/recal. & history-assisted calibration & stable history increment \\
Reuse/recal. & adaptive sentinel allocation & matched-risk incremental value \\
Selection & certified build selection & jointly feasible candidate \\
Selection & portfolio racing & Stage-I feasibility \\
Selection & protected-edge repair & native recall/budget response \\
Selection & behavioral operator search & Pareto improvement \\
Selection & critical-frontier stabilization & structure-to-budget bridge \\
Portal & fixed portal union & absolute safety and cost \\
Portal & per-query portal oracle & deployable observability \\
Portal & truth-free trigger & positive validation advantage \\
Portal & shared frontier & realizable cost compression \\
Auditor & ordered certification & native-action nesting semantics \\
Auditor & certified fallback ladder & useful intermediate fallback \\
Auditor & unified ICBA decision & decision regret and economics \\
\bottomrule
\end{tabular}
\end{table}

The ladder preserves three gaps: an oracle action may exist without an observable signal; an observable
predictor may lack certification margin; and a certifiable action may have no mean, tail, or break-even
value. Candidate attainability precedes all three.

\section{Claim Registry}
\label{app:claims}

The final theory removes several overclaims: raw \texttt{efSearch} is not assumed to generate nested
population failure events; structural similarity does not guarantee budget stability; query bootstrap is
not build-level inference; a maximum registered action is not a safe fallback without evidence; mean
savings do not imply tail savings; the negative result is transcript limited rather than an absolute
hidden-build theorem; neither main theorem is a new general Le Cam, KL, Learn-Then-Test, or best-arm
result.

\paragraph{Permitted claims.}
Registered-family HNSW evidence at 100K and the preregistered 1M replication; stage-specific supporting
Vamana evidence; fixed-target certification semantics; the deterministic same-environment contract with its
chain ablation; robust source pooling as an uncertified registered-family baseline with measured scale
dependence; grid-robustness at phenomenon level; transcript-level indistinguishability for the registered
probe class; operator-dependent cost.

\paragraph{Prohibited claims.}
Universal Graph-ANNS behavior; open-world recovery; unseen-build safety; native-budget numerical
equivalence; universal profiling or conservative-cost conclusions; production/SOTA performance;
resolution-invariant risk levels; learned-predictor transfer beyond the probed feature class.

\section{Reproducibility and Artifact Manifest}
\label{app:repro}

The final evidence derives from frozen hnswlib, clean Faiss HNSW, and DiskANN3/Vamana-style result trees,
plus the preregistered SIFT-1M cell whose manifest (design, seeds, orders, grid, roles, stop rules) was
committed before any build or search; the forensics gate passed before any risk computation. Each stage
records implementation version, build configuration, input ordering, query-role hashes, native action
grid, estimand, confidence procedure, robustness deletions, and output checksums. The clean Faiss rerun
verifies zero query/base content overlap and 48/48 frozen index hashes; the deterministic contract verifies
byte-identical indexes and search outputs at 100K and 1M. The historical 123-entry baseline remains
conditionally reproducible at 111 matching entries, eight content differences, and four unavailable
unversioned cache files. The anonymous artifact omits identifying repository links and local paths; all
revision-cycle numbers trace to committed result tables with deterministic test suites (16, 24, 27, 16, 7,
22, and 17 checks per phase). Remaining code-audit boundaries---primitive rather than complete profiling
cost, hash-only role-access status, and the three origin-unlocated historical aggregates---are recorded
explicitly and are not used to strengthen the scientific claims.

\paragraph{Reproducibility statement.}
The manuscript reports registered build families, native action grids, query-role separation, endpoint
semantics, uncertainty units, and deterministic replay checks. The anonymous artifact contains scripts,
compact result tables, manifests, and checksums without author-identifying paths. No future-replication
truth or outcome contributes to the reported analyses; any metadata or hash-only access is separately
disclosed in the artifact audit.

\paragraph{AI use statement.}
Generative AI tools assisted with organizing research artifacts, developing theory--experiment crosswalks,
checking document structure, drafting and editing manuscript text, and supporting code and artifact
validation. All mathematical statements, proofs, experimental results, citations, and AI-assisted content
require final author review, and the author assumes responsibility for the submitted manuscript.

\end{document}
